% !TeX root = ../Book1.tex
% 路径逐点相乘与先拼接再左平移的定端点同伦。
\begin{figure}[htbp]
	\centering
	\begin{tikzpicture}[font=\small,scale=1.05]
		\fill[GeometryBlueLight] (0,0) rectangle (3.8,3.2);
		\draw[GeometryGray] (0,0) rectangle (3.8,3.2);
		\draw[GeometryBlue,very thick,->] (0,0)--(3.8,3.2) node[midway,sloped,above] {\(\alpha\cdot\beta\)};
		\draw[orange!75!black,very thick,->] (0,0)--(3.8,0) node[midway,below] {\(\alpha\)};
		\draw[orange!75!black,very thick,->] (3.8,0)--(3.8,3.2) node[midway,right] {\(a\beta\)};
		\fill[GeometryBlue] (0,0) circle[radius=2pt];
		\fill[GeometryBlue] (3.8,3.2) circle[radius=2pt];
		\node[below left] at (0,0) {\((0,0)\)};
		\node[above right] at (3.8,3.2) {\((1,1)\)};
		\node at (1.9,3.85) {\(Q(s,t)=\alpha(s)\beta(t),\quad a=\alpha(1)\)};
		\draw[->,GeometryGray] (-0.8,0)--(-0.8,3.2) node[above] {\(t\)};
		\draw[->,GeometryGray] (0,-0.85)--(3.8,-0.85) node[right] {\(s\)};
	\end{tikzpicture}
	\caption{方形内的对角线与“先沿底边、再沿右边”的折线在固定端点下同伦。图中路径名称表示经 \(Q\) 映射后的路径；因此逐点乘积 \(\alpha\cdot\beta\) 与拼接 \(\alpha*(a\beta)\) 给出同一个道路类。}
	\label{g1:fig:cover-path-product}
\end{figure}
